% Source: https://artofproblemsolving.com/community/c3h2488107_pi_day__albert_einsteins_birthday_problems
% Stable identity: aops-topic-2488107-q08
If $n>m$, the left side is $0$ and the right side is positive, so there is no solution. Let $n\le m$. Then
\[\dbinom mn=\dfrac{m!}{n!\,(m-n)!}=\dfrac mn\cdot\dfrac{(m-1)!}{(n-1)!\,(m-n)!}=\dfrac mn\dbinom{m-1}{n-1},\]
so the equation is equivalent to $\dbinom{m-1}{n-1}=1$. A binomial coefficient $\dbinom{N}{K}$ with $0\le K\le N$ equals $1$ only when $K=0$ or $K=N$: for $1\le K\le N-1$ it counts the $K$-element subsets of an $N$-element set, and there are at least two of these (two different subsets of size $K$ exist because $0<K<N$). Here $n-1=0$ is excluded because $n\ne1$, so $n-1=m-1$, i.e. $m=n$. Conversely, $\dbinom nn=1=\dfrac nn$.

The solutions are $\boxed{(m,n)=(k,k)\text{ for every integer }k\ge2}$.
